\chapter{Méthodologie}
\label{chap:chap_3}
\section*{Introduction}
Sujet amené
Sujet posé
Sujet divisé


\section{Description et formulation du problème}
Présenter les sous-sections...
\subsection {Architecture proposée}
Dans ce projet, ... La figure \ref{usine} donne un aperçu de l’installation dans l’environnement RoboDK, montrant la batterie dans l’espace de travail.


\begin{figure}[!htbp]
    \centering
    \includegraphics[width=\linewidth]{usine.JPG}
    \caption{Usine de démontage d’une batterie NISSAN Leaf \cite{vanel}}
    \label{usine}
\end{figure} 




\subsection{Une autre sous-section}

\section{Une section}
Présenter les sous-sections... 



\subsection{Objectifs et contraintes}

L'équation (\ref{eq:time_matrix}) décrit...

\begin{equation}
T = \begin{pmatrix}
t_{11} & t_{12} & \cdots & t_{1n} \\
t_{21} & t_{22} & \cdots & t_{2n} \\
\vdots & \vdots & \ddots & \vdots \\
t_{n1} & t_{n2} & \cdots & t_{nn}
\end{pmatrix}
\label{eq:time_matrix}
\end{equation}

L'algorithme \ref{alg:generate_population1} présente...

\begin{minipage}{0.98\linewidth} % Réduire la largeur du cadre à 80% de la largeur de la page
\begin{algorithm}[H]
\caption{Algorithme de génération de la population initiale}
\label{alg:generate_population1}
    \KwEntree{Nombre de population, Nombre de tâches, Matrice de précédence }
  \KwSortie{Population initiale}
  Initialiser une liste vide pour stocker les chemins constituant la population \;
   \While{la taille de la population < au Nombre de population}{
    Générer un chemin aléatoire en respectant les contraintes de précédence \;
    \eIf{le chemin généré est unique et respecte les contraintes}{
     Ajouter le chemin à la population \;
     Continuer \;
     }{
     Ignorer le chemin \;
    }
   }
\end{algorithm}
\end{minipage}
\newline 
Pour améliorer la performance de l'algorithme...



\subsection{Mise en œuvre de l'algorithme génétique} 

\section{Conclusion}
Rappel des objectifs du chapitre
Indiquer l'atteinte des objectifs
Présenter le prochain chapitre

